\chapter{天を地に引きずりおろす}

\section{専攻研究}
神がおわす天界の法則を見つけた、と、そう言い始めたのはケプラーだった。
ケプラーは、天体観測に関する病的なデータを、これまた病的につぶさに検証し、ケプラーの法則を見出した。
\begin{enumerate}
  \item 惑星は、太陽を焦点の一つとする楕円を描く
  \item 惑星の単位時間の運行測度は、太陽を中心として進んだ弧が掃く面積が一定であるように進んでいる
\end{enumerate}

ケプラーまでは、天界の理論だった。
ニュートンからは違う。それはそこらの石ころが落ちるのと同じ現象だという。
ニュートンを信じてやってみよう…！



\section{モデリング}

太陽をおく。質量は$M$であるとしよう。
地球をおく。質量は$m$であるとしよう。
この二点間に、距離の逆2乗ではたらく力がかかっているとする。これが一つ目の仮定。
互いに万有引力もので引き合っている、という仮定である。
解くべき運動方程式は以下となった。
\[
  m\frac{d^2r}{dt^2} = -\frac{GMmr}{|r|^3}
\]
これを解いていこう。

